\setcounter{section}{17}

  \zwischenueberschrift{Übungsaufgaben}

\inputaufgabe
{}
{

Es sei
\maabbeledisp {\varphi} {\R^n } {\R
} { \left(  x_1   , \,   \ldots  , \,   x_n                 \right) } {a_1x_1 + \cdots + a_nx_n
} {,}
eine
\definitionsverweis {Linearform}{}{.}
Es sei $M$ der
\definitionsverweis {Graph}{}{}
dieser Funktion, den wir als
\definitionsverweis {riemannsche Mannigfaltigkeit}{}{}
auffassen. Zeige, dass zwischen den Volumina entsprechender Teilmengen des $\R^n$ und des Graphen eine konstante Beziehung besteht.

}
{} {}

\inputaufgabegibtloesung
{}
{

Berechne den Flächeninhalt der Sphäre mithilfe von
[[Graph einer Funktion/Riemannsche Untermannigfaltigkeit des R^n/Volumenform/Fakt|Korollar 17.1]].

}
{} {}

\inputaufgabe
{}
{

Diskutiere die
\definitionsverweis {Rotationsfläche}{}{} $S$ zu

\mathdisp {M= \left\{ ( \sin  { \frac{   1  }{  y } } ,y)  \mid   y> 0        \right\}} {  }

um die $x$-Achse $A$. Ist $S$ eine
\definitionsverweis {abgeschlossene Untermannigfaltigkeit}{}{} von
\mathl{\R^3 \setminus A}{?} Ist die Menge $S$
\definitionsverweis {abgeschlossen}{}{} in $\R^3$? Ist der
\definitionsverweis {Abschluss}{}{} von $S$ in $\R^3$ eine Mannigfaltigkeit?

}
{} {}

\inputaufgabegibtloesung
{}
{

Zeige, dass der Flächeninhalt der Rotationsfläche, die entsteht, wenn man den Graphen

\mathdisp {\Gamma = \left\{ (x,e^x)  \mid   x \leq 0        \right\}} {  }

um die $x$-Achse rotieren lässt, kleiner als $10$ ist.

}
{} {}

\inputaufgabe
{}
{

Bestätige, dass die in
[[Kugeloberfläche/Koordinaten von Zylinder aus/Horizontale Projektion/Flächenberechnung/Beispiel|Beispiel 17.6]],
[[Kugeloberfläche/Koordinaten von Zylinder aus/Mittelpunktsprojektion/Flächenberechnung/Beispiel|Beispiel 17.7]]
und
[[Kugeloberfläche/Geozentrische Koordinaten/Flächenberechnung/Beispiel|Beispiel 17.8]]
angegebenen
\definitionsverweis {Abbildungen}{}{}
ihr
\definitionsverweis {Bild}{}{}
auf der
\definitionsverweis {Einheitssphäre}{}{}
haben und bis auf eine
\definitionsverweis {Nullmenge}{}{}
surjektiv sind.

}
{} {}

\inputaufgabe
{}
{

Bestimme die
\zusatzklammer {partiell definierten} {} {}
\definitionsverweis {Umkehrabbildungen}{}{}
zu den in
[[Kugeloberfläche/Koordinaten von Zylinder aus/Horizontale Projektion/Flächenberechnung/Beispiel|Beispiel 17.6]],
[[Kugeloberfläche/Koordinaten von Zylinder aus/Mittelpunktsprojektion/Flächenberechnung/Beispiel|Beispiel 17.7]]
und
[[Kugeloberfläche/Geozentrische Koordinaten/Flächenberechnung/Beispiel|Beispiel 17.8]]
angegebenen
\definitionsverweis {Abbildungen}{}{.}

}
{} {}

\inputaufgabe
{}
{

Zeige, dass
\definitionsverweis {Längenkreise}{}{}
und
\definitionsverweis {Breitenkreise}{}{}
auf der Erdkugel
\definitionsverweis {senkrecht}{}{}
aufeinander stehen.

}
{} {}

\inputaufgabe
{}
{

Wie lange ist der $30$-ste
\definitionsverweis {Breitenkreis}{}{} auf der Erde
\zusatzklammer {man setze den Erdradius mit $6370$ km an} {} {.}

}
{} {}

\inputaufgabe
{}
{

Welcher Prozentanteil der Erde wird in einem Moment von der Sonne beschienen
\zusatzklammer {die Sonne soll wegen ihrer großen Entfernung als punktförmig angesetzt werden} {} {?}

}
{} {}

\inputaufgabe
{}
{

Betrachte den Ellipsoid

\mathdisp {M= \left\{ (x,y,z)\in \R^3  \mid   \ \frac{x^2}{a^2}+\frac{y^2}{b^2}+\frac{z^2}{c^2} =1         \right\}} { , }

für $a,b,c\in\mathbb{R}_{+}$. Berechne den Flächeninhalt von M für
\mathl{a= b}{.} Was passiert für
\mathl{a\rightarrow c}{?}

}
{} {}

\inputaufgabe
{}
{

Bestimme das
\definitionsverweis {Infimum}{}{}
und das
\definitionsverweis {Supremum}{}{}
der
\definitionsverweis {Länge}{}{}
der Bilder der
\definitionsverweis {Großkreise}{}{}
auf der in
[[Kugeloberfläche/Koordinaten von Zylinder aus/Horizontale Projektion/Flächenberechnung/Beispiel|Beispiel 17.6]]
beschriebenen Karte.

}
{} {}

\inputaufgabe
{}
{

Bringe
[[Graph einer Funktion/Riemannsche Untermannigfaltigkeit des R^n/Volumenform/Fakt|Korollar 17.1]]
und
[[Graph/Gradient und Volumenform/Beispiel|Beispiel 15.10]]
mithilfe von
[[Faser/Reguläre Funktion/Volumenform/Gradient und Skalarprodukt/Aufgabe|Aufgabe 16.10]]
miteinander in Verbindung.

}
{} {}

\inputaufgabe
{}
{

Es sei
\maabbdisp {\alpha} {S^2 \setminus \{N\}} { \R^2
} {}
die
\definitionsverweis {stereographische Projektion}{}{.}

a) Bestime
\mathl{\alpha_* \omega}{} zur
\definitionsverweis {kanonischen Flächenform}{}{}
auf $S^2$.

b) Berechne mit
\mathl{\alpha_* \omega}{} den Flächeninhalt der Sphäre.

}
{} {}

  \zwischenueberschrift{Aufgaben zum Abgeben}

\inputaufgabe
{5}
{

Wir betrachten den
\definitionsverweis {Graphen}{}{}
$M$ der Funktion
\maabbeledisp {\psi} {\R^2} {\R
} {(x,y)} { y+x^2
} {,}
als
\definitionsverweis {riemannsche Mannigfaltigkeit}{}{.}
Berechne den Flächeninhalt des Graphen oberhalb des Quadrats
\mathl{[-1,1] \times [-1,1]}{.}

}
{} {}

\inputaufgabe
{5}
{

Es sei

\mathdisp {M= \left\{ (x,x^2)  \mid   x \in \R        \right\}  \subseteq \R^2} {  }

die
\definitionsverweis {Parabel}{}{,}
also der
\definitionsverweis {Graph}{}{}
der
\definitionsverweis {Funktion}{}{}
\maabbeledisp {} {\R} {\R
} { x } { x^2
} {.}
Zeige, dass die zugehörige
\definitionsverweis {Rotationsfläche}{}{}
um die $x$-Achse keine
\definitionsverweis {Mannigfaltigkeit}{}{}
ist.

}
{} {}

\inputaufgabe
{4}
{

Man stelle eine
\definitionsverweis {Kugeloberfläche}{}{} als
\definitionsverweis {Rotationsfläche}{}{} dar und berechne damit den Inhalt der Kugeloberfläche.

}
{} {}

\inputaufgabe
{4}
{

Man stelle einen
\definitionsverweis {Torus}{}{} als
\definitionsverweis {Rotationsfläche}{}{} dar und berechne damit seinen Flächeninhalt.

}
{} {}

\inputaufgabe
{6}
{

Bestimme den \anfuehrung{Abstand}{}  zwischen [[w:Osnabrück|Osnabrück]] und [[w:Bangalore|Bangalore]]
\zusatzklammer {den Erdradius mit $6370$ km ansetzen} {} {}
in den beiden folgenden Sinnen.

a) Entlang der Erdoberfläche
\zusatzklammer {Luftlinie} {} {.}

b) Durch die Erde
\zusatzklammer {Maulwurfslinie} {} {.}

}
{} {}

\inputaufgabe
{6}
{

Wie lange ist das Bild des $30$-sten
\definitionsverweis {Breitenkreises}{}{}
auf den in
[[Kugeloberfläche/Koordinaten von Zylinder aus/Horizontale Projektion/Flächenberechnung/Beispiel|Beispiel 17.6]],
[[Kugeloberfläche/Koordinaten von Zylinder aus/Mittelpunktsprojektion/Flächenberechnung/Beispiel|Beispiel 17.7]]
und
[[Kugeloberfläche/Geozentrische Koordinaten/Flächenberechnung/Beispiel|Beispiel 17.8]]
beschriebenen
\definitionsverweis {Karten}{}{}
\zusatzklammer {man setze den Erdradius mit $6370$ km an} {} {?}

}
{} {}

 [[Kategorie:Latexseite]]
